\chapter{Architektur}\label{ch:architektur}

% [Cross-Ref Konsolidierte Doku, K-I.2 2026-05-15 + M.6 2026-05-18]
% Architektur-Master REV7.7: docs/architektur/02_aktueller_master_REV7_7.md
% Schichten-Modell M (4-Subsystem-Trennung): docs/architektur/10_schichten_modell_M.md
% Anti-Vermischung 4-Konzept-Ebenen: docs/architektur/11_axes_vs_strategies_disambiguation.md
% REV-Historie REV3 -> REV7.7: docs/architektur/01_REV_Historie.md
% Saeule A (Such-Strategien): docs/architektur/03_konzepte_saeule_a.md
% Saeule B (Cache-Strategien + Plattform): docs/architektur/04_konzepte_saeule_b.md
% UML pro Saeule: docs/architektur/05_uml_klassen.md
% ER-Modell: docs/architektur/06_er_modell.md
% Cross-Reference Modul x Konzept x Paper: docs/architektur/07_cross_reference.md
% Taxonomien: docs/architektur/09_taxonomien.md

Die Architektur folgt der REV-7-Skizze \cite{leis2013art}\footnote{
Der Verweis ist symbolisch --- die Architektur-Skizze REV~7 selbst liegt in
\texttt{20260508 Termin 7/Phase5\_UML\_Detail/24\_architektur\_skizze\_REV7\_2026\_05\_13.md}
und konsolidiert in \texttt{docs/architektur/02\_aktueller\_master\_REV7\_7.md}.}
und gliedert sich in drei Schichten plus Drei-Repo-Aufteilung.
Die Trennung des Cache-Engine-Werkzeugs vom Cache-Engine-Builder als
autonomes Plattform-Ausmess-System sowie die bidirektionale
Beziehung zwischen Cache-Engine und Pruefling werden in
Abschnitt~\ref{sec:vier-subsysteme} formalisiert.

\section{Vier-Subsystem-Trennung (M-Modell)}\label{sec:vier-subsysteme}

% Quelle: docs/architektur/10_schichten_modell_M.md (Phase M.2, 2026-05-18)

Diplomarbeit, Cache-Engine-Builder, Cache-Engine und Pruefling sind
\emph{vier formal getrennte Subsysteme}, die zwar im selben Projekt
wohnen, aber unterschiedliche Verantwortungen tragen. Die Trennung
ist orthogonal zur Drei-Repo-Aufteilung
(\texttt{Diplomarbeit/Code/}, \texttt{comdare-cache-engine},
\texttt{comdare-prt-art}): Cache-Engine-Builder und Cache-Engine
wohnen beide im \texttt{cache-engine}-Repo, sind dort aber zwei
unabhaengige Subsysteme (eine Executable + eine Bibliothek).

\subsection{Subsystem-Rollen}

\begin{enumerate}
    \item \textbf{Diplomarbeit/Code/messung\_driver} --- der
    Auswertungs-Orchestrator (Outer-Loop). Liest die drei
    Messreihen-XML-Konfigurationen (A/B/C), triggert pro Reihe
    den Cache-Engine-Builder und sammelt am Ende die
    Mess-Resultate fuer den Anhang.

    \item \textbf{CacheEngineBuilder} (\texttt{cache-engine/apps/}) ---
    ein \emph{autonomes Plattform-Ausmess-System}. Implementiert die
    Sieben-Phasen-Pipeline (\textsc{discover}, \textsc{measure},
    \textsc{classify}, \textsc{publish}, \textsc{bind},
    \textsc{execute}, \textsc{compare}) und enumeriert
    pro registrierter Engine den Permutations-Raum. Verwendet die
    Cache-Engine als Werkzeug, ist aber selbst eine eigenstaendige
    Executable, die ohne den \texttt{messung\_driver} aufgerufen
    werden kann.

    \item \textbf{CacheEngine} (\texttt{cache-engine/libs/cache\_engine/}) ---
    die \emph{Werkzeug-Bibliothek}. Bietet zwoelf interne
    Sub-Engines $C_1$--$C_{12}$ (Layout, Pinning, Prefetch,
    Coherence, Telemetry, Allocation, Migration, Encoding, Heuristik,
    Topologie, Scheduler, Filter), 29 Cache-Strategien $F_1$--$F_{29}$
    und etwa achtzig bis fuenfundachtzig
    \texttt{IHeuristic}-Instanzen als C++23-Concept-API.

    \item \textbf{Pruefling (PRT-ART)} (\texttt{comdare-prt-art/}) ---
    eine als \texttt{IExecutingEngine} bei der Cache-Engine
    registrierte Algorithmus-Implementation. Implementiert
    Ebene IV (Such-Engine-Familien) fuer alle dreizehn
    Algorithmus-Achsen und konsumiert die Cache-Engine-Services
    waehrend eigener Lookup-, Insert- und Scan-Operationen.
\end{enumerate}

\subsection{Bidirektionale Cache-Engine~$\leftrightarrow$~Pruefling-Beziehung}

Die symmetrische Beziehung zwischen Cache-Engine und Pruefling
ist die zentrale Eigenschaft des M-Modells:

\begin{itemize}
    \item \textbf{Richtung~(a) CE~$\rightarrow$~PRT-ART:} Die
    Cache-Engine instantiiert PRT-ART als
    \texttt{IExecutingEngine} pro Permutations-Konfiguration.
    Dies entspricht Phase~5 (\textsc{bind}) der
    Builder-Pipeline.

    \item \textbf{Richtung~(b) PRT-ART~$\rightarrow$~CE:} PRT-ART
    ruft Cache-Engine-Services (Telemetry, Prefetch,
    Heuristik\,\dots) waehrend eigener Lookups zur Laufzeit
    in Phase~6 (\textsc{execute}). Das Pruefling-Modul nutzt
    also die Cache-Engine, um seine eigene Performance
    auszumessen und zu optimieren.
\end{itemize}

Beide Richtungen muessen im Klassendiagramm (siehe U10 in
\texttt{docs/architektur/05\_uml\_klassen.md}) und im
drawio-Master-Tab (siehe Phase~P.2,
\texttt{phase5\_uml\_detail\_REV8.drawio}) explizit modelliert sein.

\subsection{Multi-Pruefling-Faehigkeit}

Die Cache-Engine kann mehrere Pruefling-Implementationen
gleichzeitig registriert halten. Messreihe~A
(\texttt{config\_a\_prt\_art\_vs\_sota.xml}) nutzt diese
Faehigkeit, um PRT-ART, ART, HOT, Masstree und weitere
SOTA-Adapter (22 V31-Adapter) in einer einzigen
Builder-Pipeline gegeneinander zu vermessen.

\section{Drei-Schichten-Hierarchie}

\begin{verbatim}
        comdare::cache_engine::CacheEngine          (Stufe 1: Stand der Technik)
                  ^
                  | erbt
        comdare::execution_engine<ProcessingStrategy>  (Stufe 2: OS-Primitiven)
                  ^
                  | erbt
        comdare::search_engine<Collection, Config>   (Stufe 3: Such-Patterns)
                  ^
                  | spezialisiert
        comdare::prt_art::PrtArtSearchEngineAdapter   (Pruefling-Anbindung)
                  |
                  | komponiert
        comdare::prt_art::PrtArtSearchEngine<Ts...>   (hybride Vector/Map/Tuple)
\end{verbatim}

\textbf{Stufe~1 (CacheEngine)} stellt die Cache-Page-Topologie, 23
Allokator-Familien, 12 Sub-Engines und Telemetrie-Strategien als
Bibliothek bereit (Stand der Technik aus den 33 Tieflektuere-Papern).

\textbf{Stufe~2 (ExecutionEngine)} wrappt diese Bausteine zu hoeheren
OS-Primitiven (cache-line-aligned Allocate, NUMA-konformes Read-Pin,
Coherence-aware Write).

\textbf{Stufe~3 (SearchEngine)} bietet Such-spezifische Patterns
(Trie-Walk, B+-Range-Scan, Hot-Path-Lookup) auf den ExecutionEngine-
Primitiven; sie ist der oberste Layer und ueberschreitet die
CacheEngine-Limits konzeptuell durch strategische Heuristiken.

\section{ABI-stabiles C++23-Modul-Interface}

Das ABI \texttt{module\_abi\_v1.hpp} folgt dem Prinzip variadischer
Template-Parameter (V12.1/V12.2): 1~Parameter $\Rightarrow$
\texttt{std::vector}-API, 2~Parameter $\Rightarrow$ \texttt{std::map}-API,
$N>2 \Rightarrow$ \texttt{std::map<K, std::tuple<V$_1$\ldots V$_N$>>}.
Komplexe Keys werden durch
\texttt{comdare::fingerprint::FixedLengthFingerprint} auf 16-Byte-
Binaer-Strings reduziert.

\section{11-Achsen-Permutationsmatrix}

Pro Stack-Ebene definiert \texttt{baustein\_variants.hpp} ein
\texttt{std::variant} ueber alle Konkretisierungen:

\begin{description}
    \item[Page] 11 Konkretisierungen (DenseByteArt256, HotMultiByte,
        MasstreeINode, CoCoCompact, Start, B2Tree, Wormhole, Surf,
        PrtArtDenseByte, PrtArtRedirect, StandardHeap)
    \item[Node] 8 (SparseNode4Art, HotBitNode, MasstreeTrie, BPlus,
        WormholeHash, SurfTrie, PrtArtRedirect, PrtArtBPlus)
    \item[Traversal] 5 (Standard, HotPath, RangeScan, RangeHotPath,
        PrtArtHotPath)
    \item[ValueHandle] 3 (Inline, External, ChainRef)
    \item[Concurrency] 6 (SingleThreaded, SharedMutex, OLC, RCU,
        HazardPointers, PrtArtOlcReserved)
    \item[Allocator] 16 (SystemMalloc, Hoard, Slab, MichaelLockFree,
        Mimalloc, Jemalloc, TCMalloc, Snmalloc, Scalloc, NUMAlloc,
        Rpmalloc, LRMalloc, CAMA, StarMalloc, Buddy, Dlmalloc)
    \item[Prefetch] 6 (None, AdaptiveDistance, PathOriented, Fractal,
        Hierarchical, PrtArtRedirectPrefetch)
    \item[Telemetry] 5 (Off, LeafOnlySampled, AllNodes, VampirOtf2,
        PikaProfile)
    \item[ISA] 6 (X86\_Baseline, SSE42, AVX2, AVX512, ARM\_Neon, ARM\_SVE)
    \item[Layout] 6 (StandardHeap, CacheObliviousVeb, CacheLineAligned,
        NumaPinned, HbmHybrid, PrtArtMultiLevel)
    \item[Reclamation] 5 (None, RcuGp, HazardPointers, EpochBased,
        CrystallineLockFree)
\end{description}

Total moegliche Permutationen: $11 \cdot 8 \cdot 5 \cdot 3 \cdot 6 \cdot
16 \cdot 6 \cdot 5 \cdot 6 \cdot 6 \cdot 5 \approx 5{,}5 \cdot 10^9$.

\section{PRT-ART als Pruefling}

Der \texttt{PrtArtSearchEngineAdapter} (V8.9, V9.1, V11.1, V13.5)
implementiert das ABI-Interface \texttt{comdare::search\_engine} via
Komposition statt direkter Vererbung. Dadurch bleibt die hybride
\texttt{PrtArtSearchEngine<Ts\ldots>}-API (66 Tests V14.1) unangetastet,
und das Adapter-Pattern erlaubt Compile-Time-Fallback auf
\texttt{cache\_engine}-Bausteine via \texttt{resolve\_baustein.hpp}
(V10.2 mit 11 Tag-Specializations).

\section{CacheEngineBuilder zwei-stufig}

\textbf{Stage~1} (CMake-Build) baut die \texttt{cache-engine-builder}-
Binary inkl. \texttt{ExperimentDriver}-Library.
\textbf{Stage~2} (Runtime, opt-in V13.4 \texttt{COMDARE\_AUTO\_RUN\_BUILDER})
ruft den Builder auf, der pro Permutation entweder ein vorkompiliertes
Modul laedt oder per \texttt{cmake/cl}-Aufruf hot-kompiliert
(\texttt{phase3\_hot\_compile\_missing}, V13.2). Optional verifiziert
\texttt{phase4b\_functional\_tests} (V13.3) den ABI-Vertrag jedes
geladenen Moduls.
